%---------------------------------------------------------------%
\addchap{LISTA DE SÍMBOLOS E ABREVIAÇÕES}
\begin{refsection}

\section*{Símbolos Estatísticos Gerais}

\begin{itemize}
\item $\Omega$: espaço amostral;
\item $P(A)$: probabilidade de ocorrência do evento $A$;
\item $E(X)$: esperança matemática da variável aleatória $X$;
\item $Var(X)$: variância da variável aleatória $X$;
\item $\sum$: operador de somatório;
\item $\varepsilon$: erro ou resíduo;
\item $\alpha$: nível de significância (probabilidade do erro tipo I);
\item $\beta$: probabilidade do erro tipo II;
\item $Z$: variável normal padronizada;
\item $t$: estatística do teste $t$ de Student;
\item $\chi^2$: estatística qui-quadrado;
\item $F$: estatística da distribuição F;
\end{itemize}

%---------------------------------------------------------------%

\section*{Parâmetros Populacionais}

\begin{itemize}
\item $\mu$: média populacional;
\item $\sigma$: desvio-padrão populacional;
\item $\sigma^2$: variância populacional;
\item $\rho$: coeficiente de correlação populacional;
\item $N$: tamanho da população;
\item $N_{h}$: tamanho da população no $h$-ésimo estrato;
\end{itemize}

%---------------------------------------------------------------%

\section*{Estatísticas Amostrais}

\begin{itemize}
\item $\bar{X}$: média aritmética amostral;
\item $\widetilde{X}$: mediana amostral;
\item $\bar{X}_{h}$: média harmônica amostral;
\item $\bar{X}_{g}$: média geométrica amostral;
\item $s$: desvio-padrão amostral;
\item $s^2$: variância amostral;
\item $r$: coeficiente de correlação amostral;
\item $CV$: coeficiente de variação;
\item $n$: tamanho da amostra;
\item $n_{h}$: tamanho da amostra no $h$-ésimo estrato;
\item $X_{i}$: valor da $i$-ésima observação da amostra;
\item $k$: número de classes em uma distribuição de frequência;
\item Valor-$p$: nível descritivo de significância do teste;
\item $H_{0}$: hipótese nula;
\item $H_{A}$: hipótese alternativa;
\end{itemize}

%---------------------------------------------------------------%

\section*{Distribuições de Probabilidade}

\begin{itemize}
\item $\lambda$: parâmetro de taxa (Poisson, Exponencial ou Gama);
\item $\theta$: parâmetro de escala (Gama);
\item $\alpha$: parâmetro de forma (Gama);
\item $\Gamma(\alpha)$: função Gama;
\item V.A.D.: variável aleatória discreta;
\item V.A.C.: variável aleatória contínua;
\end{itemize}

%---------------------------------------------------------------%

\section*{Siglas Institucionais e Técnicas}

\begin{itemize}
\item ABNT: Associação Brasileira de Normas Técnicas;
\item IBGE: Instituto Brasileiro de Geografia e Estatística;
\item DETRAN-PA: Departamento de Trânsito do Estado do Pará;
\item DATASUS: Departamento de Informática do Sistema Único de Saúde do Brasil;
\item FAPESPA: Fundação Amazônia de Amparo a Estudos e Pesquisas;
\item RAIS: Relação Anual de Informações Sociais;
\item FINBRA: Finanças do Brasil (Secretaria do Tesouro Nacional);
\item SPSS: Statistical Package for the Social Sciences;
\item SAS: Statistical Analysis System;
\item RStudio: Ambiente de desenvolvimento integrado para a linguagem R;
\end{itemize}

\printbibliography[heading=subbibliography]
\end{refsection}
%---------------------------------------------------------------%